%------------------------------------------------------------------------------
% \PART 3
%------------------------------------------------------------------------------
\chapter[Разработка нейронной сети для распознавания рукописных цифр]
{РАЗРАБОТКА НЕЙРОННОЙ СЕТИ ДЛЯ РАСПОЗНАВАНИЯ РУКОПИСНЫХ ЦИФР}

%------------------------------------------------------------------------------
% \PART 3.1 Text
%------------------------------------------------------------------------------
\section{Функциональная спецификация системы}\par
\hspace*{12.5 mm}Для разработки системы распознавания рукописных цифр с 
использованием нейронной сети прямого распространения необходимо составить 
функциональную спецификацию, которая будет описывать требования к системе, её 
функциональные компоненты и интерфейсы с пользователем и окружением.

Система будет реализована как IP-блок нейронной сети прямого распространения, 
которая предварительно будет обучаться на наборе данных изображений рукописных 
цифр, а затем использоваться для их распознавания.

В контексте требований пользователя система должна отвечать на следующие 
вопросы:

  1 Какие средства необходимо предусмотреть для ввода данных?

  2 Какие средства необходимо предусмотреть для вывода результатов?

  3 Как будет осуществляться хранение весовых коэффициентов?

  4 Какие средства необходимы для работы с пользователем?

Составим функциональную спецификацию для разрабатываемой системы:

  1 Данные о рукописных цифрах будут передаваться в виде последовательной 
передачи пикселей изображений, которые будут поступать в IP-блок по AXI4-lite 
интерфейсу из процессорной системы на ПЛИС.
    
  2 Результаты работы модели будут поступать в процессорную систему по 
интерфейсу AXI4-lite, где проходить обработку, например использоватья в 
построении матрицы ошибок.
    
  3 Хранение весовых коэффициентов будет осуществляться в блочной памяти. 
Необходимо обеспечить поддержку трех весовых групп: первого слоя LST, 
второго слоя LST и выходного полносвязного слоя.
    
  4 Система будет предоставлять IP-блок, который будет подключаться по 
AXI4-lite интерфейсу к процессорной системе.

%------------------------------------------------------------------------------
% \PART 3.2
%------------------------------------------------------------------------------
\section{Разбиение системы на модули}\par
\hspace*{12.5 mm}Система распознавания рукописных цифр на базе нейронной сети 
прямого распространения может быть разделена на несколько функциональных 
модулей:

  1 Модуль памяти изображения, который хранит изображение на всех этапах.

  2 Модуль обработки пикселя, который передает пиксель и необходимый весовой 
коэффициент в модуль MAC.\@
  
  3 Модуль MAC выполняет операцию умножения с накоплением.

  4 Модули памяти весов, в которых хранятся весовые коэффициенты.

  5 Модуль функции активации тангенс гиперболический, который выполняет 
вычисление функции тангенс гиперболический в LST слое.

  6 Модуль функции активации softmax.

  7 Модуль управления.

%------------------------------------------------------------------------------
% \PART 3.3
%------------------------------------------------------------------------------
\section{Принцип работы LST преобразования}\par
\hspace*{12.5 mm}Двумерное разделимое преобразование используется в обработке 
изображений для уменьшения вычислительной сложности. Основная идея заключается 
в том, что вместо хранения полного двумерного ядра свёртки размером 
n\(\times\)n используется его представление в виде произведения двух одномерных 
векторов. Это позволяет сократить число параметров с \(n^2\) до \(2n\), что 
снижает затраты на вычисления и память.

\begin{equation}
    W = V \times h^T,
\end{equation}

\noindentгде \(W\in \mathbb{R}^{n\times n}\);\\ 
\hphantom{где} \(v, h^T\in \mathbb{R}^{n\times 1}\).

LST основан на идее совместного использования весов одного полносвязанного 
слоя для обработки всех строк изображения. После этого второй общий 
полносвязный слой используется для обработки всех столбцов представления 
изображения, полученных из первого слоя. Использование слоев LST в архитектуре 
нейронной сети значительно сокращает количество параметров модели по сравнению 
с моделями, которые используют стекированные полносвязные слои.

LST можно использовать для построения архитектур глубоких нейронных сетей 
(DNN), заменяя традиционную нейронную сеть прямого распространения. 

Предложенное обученное двумерное разделяемое преобразование
обрабатывает изображение построчно, а не преобразует его в вектор, как это 
обычно делается в полносвязных сетях. 

С математической точки зрения LST2D принимает двумерное изображение X размером 
din\(\times\)din в качестве входных данных и создает двумерное выходное 
изображение Y размером dout\(\times\)dout

\begin{equation}
  Y = LST_{d_{in}xd_{out}}(X) = \tan(W_2\tan(W_1X^T)).
\end{equation}

\noindentгде W1, W2 — матрицы весов слоев FC1 и FC2 соответственно.

Принцип работы LST2D представлен на рисунке~\ref{fig:lst2d-pic}\cite{LST2D}. 

\insertfigure{lst2d-pic}{pic/LST2D.png}{Принцип работы LST2D}{16cm}

Здесь din и dout — гиперпараметры преобразования, определяющие количество 
обучаемых параметров N=2 $\cdot$ (din+1) $\cdot$ dout = 2 $\cdot$ 29 $\cdot$ 28 = 1624. 

Простейшая архитектура нейронной сети на основе LST2D для распознавания цифр 
MNIST требует одного блока LST2D, за которым следует полносвязный слой с 
функцией активации softmax, что показано на 
рисунке~\ref{fig:lst-pic}\cite{LST2D}. 

\insertfigure{lst-pic}{pic/LST.png}{Принцип работы LST-1}{10cm}

Алгоритм работы LST-1 преобразования представлен на рисунке~\ref{fig:alg-lst}.
\insertfigure{alg-lst}{pic/alg_lst.png}{Алгоритм работы LST2D слоя}{16cm}

Данную архитектуру можно масштабировать. Пример двухблочной архитектуры LST2D 
приведен на рисунке~\ref{fig:lst2-pic}\cite{LST2D}.

\insertfigure{lst2-pic}{pic/LST2.png}{Принцип работы LST-2}{11cm}

Пример обработки изображения нейронной сетью архитектуры LST-1 представлен на 
рисунке~\ref{fig:img-proc}\cite{LST2D}.

\insertfigure{img-proc}{pic/img_proc.png}{Архитектура CNN-Transformer}{16cm}

%------------------------------------------------------------------------------
% \PART 3.4
%------------------------------------------------------------------------------
\section{Обучение нейронной сети}
\hspace*{12.5 mm}Обучение нейронной сети LST-1 на основе 60000 тренировочных 
изображений из базы данных MNIST.\@ Для обучения используется фреймворк 
PyTorch.\@

В качестве исходных данных используется набор MNIST, 
содержащий 70 тысяч полутоновых изображений размера $28 \times 28$ пикселей, 
представляющих рукописные цифры от 0 до 9~\cite{MNIST}. Данный набор разделен 
на две части:

  – \text{Тренировочная выборка} – 60 тысяч изображений.

  – \text{Тестовая выборка} – 10 тысяч изображений.

    Пример данных из базы MNIST представлен на рисунке~\ref{fig:mnist}.

\insertfigure{mnist}{pic/MNIST.jpeg}{Пример данных из MNIST}{11cm}

Из тренировочного набора 
выделяется 1000 изображений для валидации, оставшиеся 59000 используются для 
непосредственного обучения модели. 

Перед подачей в нейросеть изображения проходят предварительную обработку, 
включающую нормализацию значений пикселей. Для упрощения 
процесса обучения выполняется нормализация данных таким образом, чтобы значение 
каждого пикселя было в диапазоне от $-1$ до $1$, а стандартное отклонение 
($\sigma$) составляло $0.5$. Это позволяет сделать градиентный спуск более 
устойчивым и ускорить процесс сходимости модели. Пример такой обработки 
представлен на рисунке~\ref{fig:norm-img}.

\insertfigure{norm-img}{pic/norm_img.png}{Пример обработанных данных из MNIST}{9cm}

Нормализация изображений осуществляется в соответствии с описанными 
требованиями:

\small
\fontsize{12pt}{12pt}\selectfont
\begin{verbatim}
    transform = transforms.Compose([
        transforms.ToTensor(),
        transforms.Normalize((0.5,), (0.5,))
    ])
\end{verbatim}
\normalsize
\fontsize{14pt}{14pt}\selectfont

Модель L2DST состоит из двух линейных слоев LazyLinear с функцией активации 
tanh и полносвязный слой Linear с функцией активации softmax для классификации. 

В качестве метрики, для оценки обучаемости нейронной сети используется функция 
потерь кросс-энтропии – {Negative Log-Likelihood Loss (NLLLoss)}.

Оптимизатор {Adam} с параметром скорости обучения \(1.5 \times 10^{-3}\) и 
шаговым снижением {StepLR} (с понижением на 3\% каждые 10 эпох).
 
\small
\fontsize{12pt}{12pt}\selectfont
\begin{verbatim}
    class L2DST(nn.Module):
    def __init__(self, ffn_num_hiddens, ffn_num_outputs):
        super().__init__()
        self.dense1 = nn.LazyLinear(ffn_num_hiddens)
        self.dense2 = nn.LazyLinear(ffn_num_outputs)

criterion = nn.NLLLoss() 
optimizer = optim.Adam(model.parameters(), lr=15e-4, weight_decay=0e-5)
scheduler = lr_scheduler.StepLR(optimizer, step_size=10, gamma=0.97, 
                                                         verbose=False)
\end{verbatim}
\normalsize
\fontsize{14pt}{14pt}\selectfont

На рисунке~\ref{fig:loss} показан график функции потерь за 370 эпохах.

\insertfigure{loss}{pic/loss.jpg}{Графики функции потерь на обучающем и тестовом наборе}{10cm}

После завершения каждой эпохи вычисляются средние значения ошибки на 
тренировочном и валидационном наборах данных.  

Python описание процесса обучения нейронной сети прямого распространения 
приведено в приложении Б.